%% Harbin Institute of Technology thesis starter — hithesis document class.
%%
%% hithesis is NOT included in this download: it ships with TeX Live (and with
%% latextemplates.com's compiler), so bundling a copy would only pin you to an
%% older class than your own distribution already has.
%%
%% Class documentation: texdoc hithesis  |  https://ctan.org/pkg/hithesis
%% Licence: hithesis.cls is LPPL-1.3a, (c) 2017-2022 Chu Yanshuo.
%%
%% XELATEX IS REQUIRED, not a preference. hithesis loads ctexbook, which loads
%% xeCJK; neither works under pdflatex, so a pdflatex run dies at the class
%% rather than producing a worse-looking PDF. See engine.txt.
%%
%% `fontset=fandol' is passed EXPLICITLY. Left unset, ctex probes the system
%% for a CJK font family and prefers Windows fonts (SimSun/SimHei) that exist
%% on almost no Linux box or CI runner; the failure is a font-not-found abort
%% at load time, on a machine where the class itself is installed correctly.
%% Fandol ships with TeX Live itself, so naming it is what makes this starter
%% compile the same way everywhere. Drop the option on a machine with the
%% official Chinese fonts if you want the exact submission typefaces.
%%
%% `type=master' selects the master's-degree title pages and single-sided
%% layout. The class also accepts type=bachelor and type=doctor, which change
%% the cover sheets, the page-numbering rules, and (for doctor) switch to
%% two-sided printing — so this is the first option to check against your own
%% degree, not a cosmetic default.
%%
%% `newtxmath=false' works around a genuine clash INSIDE the class, not a
%% preference. hithesis loads newtxmath and then, ~80 lines later, amssymb;
%% both define \Bbbk, so the default settings abort with
%%   ! LaTeX Error: Command `\Bbbk' already defined.
%% before a single page is typeset. The class ships this boolean as its own
%% escape hatch. Turning it off gives you the Times math from newtxtext
%% rather than newtxmath's; if you need newtxmath specifically, load amssymb
%% yourself before the class instead.
\documentclass[type=master,fontset=fandol,newtxmath=false]{hithesis}

%% Every cover-sheet field is a key of \hitsetup. They are separate Chinese and
%% English keys rather than one bilingual field because hithesis typesets TWO
%% covers: the Chinese title pages and, for master's/doctoral theses, a
%% separate English cover. Leaving an e... key empty does not fall back to the
%% Chinese one; it prints blank on the English cover.
\hitsetup{
  ctitle       = {哈尔滨工业大学学位论文模板},
  etitle       = {A Thesis Template for Harbin Institute of Technology},
  cauthor      = {张三},
  eauthor      = {Zhang San},
  cstudentid   = {21S000000},
  cdate        = {2026年8月},
  edate        = {August, 2026},
  cxueke       = {工学},
  exueke       = {Engineering},
  csubject     = {计算机科学与技术},
  esubject     = {Computer Science and Technology},
  caffil       = {计算机科学与技术学院},
  eaffil       = {School of Computer Science and Technology},
  csupervisor  = {李四教授},
  esupervisor  = {Prof. Li Si},
  %% Keywords are COMMA-SEPARATED LISTS, not free text: the class splits on the
  %% comma and re-joins with its own separator (a Chinese enumeration comma for
  %% ckeywords, "; " for ekeywords). Typing your own separator here gets it
  %% printed literally alongside the class's.
  ckeywords    = {学位论文, 模板, 排版},
  ekeywords    = {thesis, template, typesetting},
}

\begin{document}

%% The two abstracts are declared BEFORE \makecover and stored, not typeset in
%% place: \makecover emits the covers and then the abstract pages in the order
%% the university requires. Moving these below \makecover silently produces a
%% thesis with empty abstract pages, because the macro has already expanded.
\begin{cabstract}
  中文摘要应当概括论文的研究背景、所采用的方法、主要结果与结论。摘要是独立
  成篇的短文，不应出现图表、公式编号或参考文献引用。
\end{cabstract}

\begin{eabstract}
  The English abstract states the problem, the method, the principal results
  and the conclusions. It is self-contained: no citations, no numbered
  equations, and no references to figures or tables in the body.
\end{eabstract}

%% \makecover, NOT \maketitle. hithesis does `\let\title\@gobble' and never
%% redefines \maketitle, so calling \maketitle here reaches LaTeX's stock one,
%% which has no title to print and stops with "! LaTeX Error: No \title given."
%% \makecover is what emits the Chinese covers, the English cover, and then
%% both abstract pages.
\makecover

\frontmatter
\tableofcontents

\mainmatter

\chapter{绪论}
\label{chap:intro}
介绍研究背景与意义，说明国内外研究现状，并给出本文的主要工作与组织结构。

\section{研究背景}
在这里展开背景论述。

\section{本文主要工作}
概述各章内容。

\chapter{相关理论与方法}
\label{chap:method}
阐述论文所依据的理论基础与所采用的方法。

\section{表格示例}
表~\ref{tab:example} 给出了一个三线表示例。

\begin{table}[htbp]
  \centering
  \caption{实验结果对比}
  \label{tab:example}
  \begin{tabular}{lrr}
    \hline
    方法 & 准确率 & 耗时/ms \\
    \hline
    基线方法 & 0.912 & 24 \\
    本文方法 & 0.947 & 18 \\
    \hline
  \end{tabular}
\end{table}

\section{公式示例}
\begin{equation}
  \label{eq:example}
  y = \sum_{i=1}^{n} w_i x_i + b
\end{equation}

式~\eqref{eq:example} 定义了模型的线性部分。

\chapter{实验与分析}
\label{chap:exp}
给出实验设置、评价指标与结果分析，并与第~\ref{chap:method} 章的方法相对应。

%% \conclusions is an ENVIRONMENT, not a \chapter: it emits the unnumbered
%% "结论" heading the university mandates, with the matching English entry in
%% the bilingual table of contents. A \chapter*{结论} would look identical on
%% the page and be missing from both contents lists.
\begin{conclusions}
  总结全文的主要结论，说明工作的创新点与不足，并展望后续研究方向。
\end{conclusions}

\begin{acknowledgements}
  感谢导师的悉心指导，感谢课题组同学与家人的支持。
\end{acknowledgements}

\end{document}
